\chapter{Other I/O classes, functions and templates}

\section{\texttt{Dwm::EncodedUnsigned<T>}}

\subsection{\texttt{Dwm::EncodedU64}}

\subsection{\texttt{Dwm::EncodedU32}}

\subsection{\texttt{Dwm::EncodedU16}}

\section{\texttt{Dwm::EncodedSigned<T>}}

\subsection{\texttt{Dwm::EncodedS64}}

\subsection{\texttt{Dwm::EncodedS32}}

\subsection{\texttt{Dwm::EncodedS16}}

\section{\texttt{Dwm::EncodedU64}}
\texttt{Dwm::EncodedU64} represents a \texttt{uint64\_t} value, but
on the wire is encoded as an encoding byte followed by a variable
length (1 to 8) sequence of bytes containing the value.  It is a
typical space/time tradeoff: it requires more time to read or write,
but saves space on the wire and storage for some common use cases.

Common use cases generally revolve around \texttt{std::size\_t} values
representing the length of a \texttt{std::string}, the number of elements
in an STL container (\texttt{vector}, \texttt{map}, et. al.), or the
return of the \texttt{sizeof} operator.  On a 64-bit platform, you may
need to accomodate 64-bit length values, but find that all of your values
are less than $2^{32}$.  Or less than 256.  These are scenarios where
\texttt{Dwm::EncodedU64} may be useful if the space/time tradeoff is
desired.

See \hyperref[encoding:DwmEncodedU64]{\texttt{Dwm::EncodedU64} on-the-wire
encoding} for details on the encoding format.

\section{\texttt{Dwm::IOUtils}}

\section{Endian conversions}
\subsection{\texttt{template <typename T>\\
    concept IsEndianSensitiveInteger}}
True if \texttt{std::is\_integral\_v<T>} and \texttt{sizeof(T)} is 2,
4 or 8.

Examples:
\begin{itemize}[noitemsep,nolistsep]
\item[] \texttt{IsEndianSensitiveInteger<uint64\_t>} $\rightarrow$ \texttt{true}
\item[] \texttt{IsEndianSensitiveInteger<uint32\_t>} $\rightarrow$ \texttt{true}
\item[] \texttt{IsEndianSensitiveInteger<uint16\_t>} $\rightarrow$ \texttt{true}
\item[] \texttt{IsEndianSensitiveInteger<uint8\_t>} $\rightarrow$ \texttt{false}
\item[] \texttt{IsEndianSensitiveInteger<char>} $\rightarrow$ \texttt{false}
\item[] \texttt{IsEndianSensitiveInteger<float>} $\rightarrow$ \texttt{false}
\end{itemize}

\subsection{\texttt{template <typename T>\\
    requires IsEndianSensitiveInteger<T>\\
    inline auto BE2Host(T t) $\rightarrow$ T}}
Returns a \texttt{T} holding the value of \texttt{t} converted from
big endian (MSB first) byte order to host byte order.  Example:
\begin{verbatim}
  uint16_t  x = 0xFF00;
  uint16_t  y = Dwm::BE2Host(x);
  // on a little endian machine, y == 0x00FF
  assert(0x00FF == y);
\end{verbatim}

\subsection{\texttt{template <typename T>\\
    requires IsEndianSensitiveInteger<T>\\
    inline auto Host2BE(T t) $\rightarrow$ T}}
Returns a \texttt{T} holding the value of \texttt{t} converted from
host byte order to big-endian (MSB first) byte order.  Example:
\begin{verbatim}
  uint32_t  x = 0xAABBCCDD;
  uint32_t  y = Dwm::Host2BE(x);
  // on a little endian machine, y == 0xDDCCBBAA
  assert(0xDDCCBBAA == y);
\end{verbatim}

\subsection{\texttt{template <typename ...Args>\\
    requires (IsEndianSensitiveInteger<Args> and ...)\\
    void BE2Host(Args \& ...args)}}

\subsection{\texttt{template <typename ...Args>\\
    requires (IsEndianSensitiveInteger<Args> and ...)\\
    void Host2BE(Args \& ...args)}}
